\documentclass[../../../uem_phy_std_q.tex]{subfiles}

\begin{document}
    $xy$平面上の電位分布を測定し，2\:V間隔で等電位線を描いたところ，%
    図のようになった．%
    等電位線の横に添えた数値は，各等電位線上での電位の値である．

    \begin{enumerate}        
        \item%
            陽子（電荷$e=1.6 \times 10^{-19}$\:C）を，%
            クーロン力に逆らってゆっくりと運ぶとする．%
            点Pから点Fへと運ぶ際に必要な仕事$W_\mathrm{PF}$，%
            点Pから点Gへと運ぶ際に必要な仕事$W_\mathrm{PG}$，%
            点Pから点Hへと運ぶ際に必要な仕事$W_\mathrm{PH}$を，それぞれ求めよ．
        \item%
            電子（電荷$-e=-1.6 \times 10^{-19}$\:C）を，%
            クーロン力に逆らってゆっくりと運ぶとする．%
            点Pから点Fへと運ぶ際に必要な仕事$W^*_\mathrm{PF}$，%
            点Pから点Gへと運ぶ際に必要な仕事$W^*_\mathrm{PG}$，%
            点Pから点Hへと運ぶ際に必要な仕事$W^*_\mathrm{PH}$を，それぞれ求めよ．
        \item 点Pを通る電気力線の概略を描け．
        \item%
            点P付近に図示してある等電位線の間隔は大体0.53\:m程度である．%
            このことから点Pにおける電場の強さ$E_\mathrm{P}$の近似値を求めよ．
    \end{enumerate}
    \centering
        \includegraphics{\subfix{fig_uem_phy_std_em_ef_04/fig_uem_phy_std_em_ef_04_03_q.pdf}}
    
\end{document}




